% Modernised reading — fonds Alexandre Grothendieck, shelfmark 58, pages 1-66
% (the whole folder, four batches of transcription: 1-20, 21-40, 41-60, 61-66).
%
% This is an INTERPRETATION, not a transcription. It re-reads the pages in
% today's notation and vocabulary, states the hypotheses the manuscript leaves
% implicit, and drops the critical apparatus so that the mathematics reads
% straight through. Everything the brackets and underlines carried — a doubtful
% reading, a notation he used differently, a missing passage — moves into a
% footnote. It is held to being mathematically correct as it stands; where the
% pages are loose or mistaken, the true statement is what appears here and the
% footnote says what the page has. In any dispute of substance the
% transcription governs: whoever cites Grothendieck must cite batch-NN.fr.tex.
%
% Scope: all four batches transcribed, read whole before any of this was
% written; ONE file for the shelfmark. The fast pages (1-20, 22-26, 29-30, 45,
% 49) are largely illegible in their prose: the reading restates what their
% formulas and statements carry and says, section by section, where it
% reconstructs and where it does not.
%
% What the folder is. Four stretches separated by three covers (pp. 1, 27,
% 53): (I) pp. 2-26, Grothendieck's connectedness theorem — local form on the
% punctured spectrum (Lemmes 2-3, Théorème, Corollaires, pi_1), local-to-global
% (Lemme 4), special fibres of a proper morphism (Th. 2, Lemme 5, Cor. 1-3);
% (II) pp. 29-51, « couples de Lefschetz-Grauert »: introduction and nine-item
% plan, the conditions (LG)_n, LG, LG stricte (today Lef, Leff of SGA 2 X),
% consequences for pi_0, pi_1, étale covers, Pic, a permanence theorem under
% projective flat morphisms, a local form over (A, t); (III) pp. 55-61, P^r,
% a hyperplane section Y, a neighbourhood U, Z = X - U finite: chains of
% equivalent conditions (local cohomology, Serre and local duality, depth);
% (IV) pp. 64-66, scratch on Pic of cones, ending mid-sentence.
%
% Notation fixed for the whole document. A local noetherian, m its maximal
% ideal (his « rA »), X = Spec A, X' = X - {m} the punctured spectrum (ONLY
% this: the scheme projective over X of pp. 41-43 is renamed W, the blown-up
% cone of p. 66 is renamed P, the base of Th. 2 is renamed B — the page calls
% them X', X', Y). Dimensions in X' are those of X' as a space. « Connexe en
% dimension >= k » as in Hartshorne 1962 (defined in the conventions; the
% pages never define it, nor the conditions a_j, b_j, c_j). The ground field
% is \Bbbk, k being the integer of connectedness. X_{/Y} = \hat X the formal
% completion. Part III: S = \Bbbk[T_0..T_r], E^j = Ext^j(F, omega_X), g : U -> X.
%
% Substantive departures from the pages, each footnoted where it occurs:
% — p. 8: the connectedness hypothesis is put on the completion (or A taken
%   complete), as in SGA 2 XIII; the page puts it on X'.
% — p. 2: Lemme 2 needs X' connected in dimension >= 1 (counterexample: two
%   3-planes meeting at a point); supplied, as the case k = 1 of p. 10.
% — p. 18: the pi_1 corollary is stated for A complete normal; the page's
%   adjective is illegible and the proof uses irreducibility of covers.
% — p. 20: Lemme 4 needs X connected; supplied. The local hypothesis needed
%   is weaker than the page's; both stated.
% — p. 22: Th. 2 needs X_1 GEOMETRICALLY connected (a number-field example for
%   k = 0); the dimension bound read k+1 (the page's « k » is doubtful). The
%   reading does not claim to have checked the proof beyond a sketch.
% — p. 33: Remarque 1's « Oui » (injectivity in degree n+1 is automatic) is
%   proved here for X projective over an affine noetherian base (ours).
% — p. 36: the fourth member of the « conclusion pratique » is an equivalence
%   (ours; his curved arrows say as much).
% — p. 40: « pleinement fidèle » / « équivalence » for invertible sheaves on
%   \hat X -> Y overstate: only the statements on Pic hold as written.
% — p. 41: « plat sur Y » read « plat sur X ».
% — p. 45: A must be t-adically complete (supplied, the page's second line is
%   overwritten); without a depth hypothesis rho is NOT full (example given).
% — p. 55: the depth bound left blank is supplied (prof F_x > n at closed points).
% — p. 57: the third member of the alpha_i-bijective group corrected; in B_n
%   the last equivalence holds only as an implication (example: a skyscraper);
%   the « (*) » before the comparison theorem read as B_n.
% — p. 64: « Spec » read « Proj ».
%
% Announced and not established: Lemme 1 and Corollaire 2 (cited pp. 4, 6, 24,
% not in the folder); conditions a_j, b_j, c_j (p. 8); the converse Corollaire I
% (pp. 12-16, largely cancelled); Cor. 1-2 to Th. 2 (one line each); the
% chapter whose § 4-5 and N° 1-2 the introduction cites; items 1-4 and 8-9 of
% the plan; (iii) and Cor. 2 of the permanence theorem (no proof); the « cor.
% 3 au th. de comparaison affine », conditions a), b), c'), set T_1 (p. 49);
% the end of the local theorem's proof (the « ? » of p. 47); the questions
% b) and c) of p. 55 beyond the sufficient conditions; Pic(C_X) = Pic(X) (p. 66,
% breaks off). Typed leaves (pp. 3-19 odd, 21-25 odd, 46-52 even, 56, 60, 62)
% are other texts with nothing in his hand but pp. 56 and 60.
%
% Pass: Opus 5.5 (claude-opus-5-5), 2026-10-04 — first pass, unchecked against the pages by a human.
\documentclass[11pt,a4paper]{article}
\input{../preamble/grothendieck}

\folder{58}
\pages{1}{66}
\dating{[à partir de 1964]}
\watermark{Édition de démonstration\\interprétation personnelle de l'œuvre}
\foldertitle{Théorie de " Lefschetz-Grauert ". Applications en π₁, à Pic etc… : notes manuscrites (s.d.) — lecture modernisée du dossier entier}

\begin{document}

\section*{Résumé}

\begin{resume}
Ces soixante-six pages tournent autour d'une seule question : que garde une
partie d'un espace de la forme de l'espace tout entier ? Le modèle est le
théorème de Lefschetz, vieux de quarante ans quand ces notes sont écrites :
une section hyperplane d'une variété projective lisse de dimension $n$ a,
jusqu'en dimension $n-2$, la même topologie que la variété. Une surface
coupée par un plan général donne une courbe connexe ; une variété de
dimension $3$ a le même groupe fondamental que ses sections par des plans
généraux. Grothendieck veut des énoncés de ce genre en géométrie algébrique
pure, sur n'importe quel corps et même sur des anneaux locaux, là où il n'y a
plus de topologie classique pour les démontrer.

La première partie traite de la connexité. Une feuille de papier ne se coupe
pas en deux en lui retirant un point ; deux feuilles collées par un seul point
le peuvent. On dit, depuis Hartshorne, qu'un espace est \emph{connexe en
dimension $\geq k$} quand on ne le déconnecte pas en lui retirant une partie de
dimension $< k$. Le théorème démontré ici dit qu'en coupant un espace connexe
en dimension $\geq k$ par $l$ équations on perd au plus $l$ crans : ce qui
reste est connexe en dimension $\geq k-l$. Il est d'abord établi au voisinage
d'un point, puis transporté à des familles : si les fibres générales d'une
famille propre sont bien connexes, la fibre spéciale l'est aussi. On en tire
la surjectivité de l'application entre groupes fondamentaux.

La deuxième partie change d'instrument. Plutôt que de couper, on regarde la
variété $X$ dans un voisinage infiniment petit d'une sous-variété $Y$ --- son
\emph{complété formel} le long de $Y$, qui retient de $X$ tous les
développements de Taylor dans les directions transverses à $Y$ et rien
d'autre. La question devient : que perd-on en ne connaissant $X$ que par ce
voisinage ? Grothendieck appelle « de Lefschetz-Grauert » les couples
$(X, Y)$ pour lesquels on ne perd rien de certaines données --- les sections
des fibrés vectoriels, puis les fibrés eux-mêmes. Il en déduit
mécaniquement des énoncés de Lefschetz pour les composantes connexes, pour le
groupe fondamental, pour le groupe de Picard des fibrés en droites, et il
montre que la propriété se transporte par les morphismes projectifs et plats.
Une troisième partie traduit tout cela, pour l'espace projectif, en conditions
de \emph{profondeur} : une mesure algébrique de la place qu'un faisceau
laisse autour de chaque point, que la dualité de Serre permet de lire sur la
cohomologie. Le dossier s'achève sur des brouillons à propos du groupe de
Picard d'un cône.

Ces pages sont le laboratoire de ce qui deviendra le séminaire SGA 2 :
les conditions qu'il nomme ici « de Lefschetz-Grauert » y portent les noms
$\mathrm{Lef}$ et $\mathrm{Leff}$. On y voit l'écriture d'un travail en cours
plus qu'un texte : un énoncé encadré puis rayé, une condition jugée «
inutile » dans la marge, « je m'aperçois que », deux points d'interrogation
devant $\operatorname{Pic}(C) = 0$, et, au bout d'une démonstration difficile,
« on a gagné ». Les questions restent écrites comme des questions, et la
dernière page s'arrête au milieu d'une phrase.

\keywords{Grothendieck connectedness theorem, connectedness in dimension k, punctured spectrum, Hartshorne connectedness theorem, Zariski connectedness theorem, Lefschetz hyperplane theorem, formal completion, Lefschetz condition, effective Lefschetz condition, étale fundamental group, Picard group, local cohomology, local duality, Serre duality, depth, finiteness theorem, parafactorial ring, projecting cone, homotopy invariance of Pic}
\end{resume}

\section*{Le fil du dossier, et les conventions}

\begin{enumerate}
\item[I.] \emph{La connexité} (pages 1 à 26, sous la couverture « Applicat.\
à $\pi$ »\footnote{Titre de sa main sur la couverture (page 1) ; le dernier
signe, lu $\pi$, est douteux et pourrait être $\pi_1$.}) : la connexité en
dimension $\geq k$ du spectre épointé d'un anneau local et de ses sections
(pages 2 à 16) ; le groupe fondamental (page 18) ; le passage du local au
global (page 20) ; les fibres spéciales d'un morphisme propre et les sections
hyperplanes projectives (pages 22 à 26).
\item[II.] \emph{Les couples de Lefschetz-Grauert} (pages 27 à 51, sous la
couverture « Couple de Lefschetz-Grauert $(X, Y)$ ») : introduction et plan
(pages 29 à 31) ; définitions et premières remarques (pages 32 à 36) ;
$\pi_0$, $\pi_1$, revêtements étales, $\operatorname{Pic}$ (pages 37 à 40) ;
permanence par morphisme projectif et plat (pages 40 à 43) ; une forme locale
(pages 45 à 51).
\item[III.] \emph{Morphismes projectifs et propres} (pages 53 à 61, sous la
couverture « Applications aux morphismes projectifs et aux morphismes propres
») : sur $\mathbf{P}^r$, un voisinage $U$ d'une section hyperplane et son
complémentaire fini $Z$ ; chaînes de conditions équivalentes.
\item[IV.] \emph{Le groupe de Picard des cônes} (pages 64 à 66) : brouillons.
\end{enumerate}

Les feuillets dactylographiés que le dossier intercale --- au dos des pages 2 à
20, les pages 21 à 25, 46 à 52 et 62 --- appartiennent à d'autres textes et ne
portent rien de sa main ; les pages 28, 44 et 54 sont blanches, la page 63 ne
porte que le décalque de la page 60. Seuls les feuillets tapés des pages 56
et 60 portent des calculs de lui.

\emph{Le fil.} Les quatre parties ne forment pas une démonstration continue
mais un programme : le plan de la page 31 l'énonce en neuf points, dont la
partie I réalise le point 8 pour la connexité (« théorèmes de Lefschetz »
sans complété formel), la partie II les points 5 à 7, la partie III un
morceau du point 4 (« cas particulier : schémas projectifs »), et la partie IV
amorce le point 9. Ce plan suit de près l'ordre des exposés VIII à XIII de
SGA~2.\footnote{A.~Grothendieck, \emph{Cohomologie locale des faisceaux
cohérents et théorèmes de Lefschetz locaux et globaux} (SGA~2), séminaire de
1962, publié en 1968 : VIII, théorème de finitude ; IX, géométrie algébrique
et géométrie formelle ; X, groupe fondamental ; XI, groupe de Picard ; XII,
schémas projectifs ; XIII, problèmes et conjectures, dont le théorème de
connexité. Le rapprochement est le nôtre ; rien sur les pages ne dit si elles
précèdent ou suivent le séminaire, et la datation de Montpellier, « à partir
de 1964 », ne tranche pas.}

\emph{Conventions.} Pour un anneau local noethérien $A$ d'idéal maximal
$\mathfrak{m}$, on note $X = \operatorname{Spec} A$ et
$X' = X \smallsetminus \lbrace \mathfrak{m} \rbrace$ son \emph{spectre
épointé}.\footnote{Il écrit $\mathfrak{r}A$ (le radical) pour $\mathfrak{m}$.}
La lettre $X'$ ne désigne que cela dans tout le document : le schéma
projectif sur $X$ des pages 41 à 43 est ici noté $W$, l'éclaté du cône de la
page 66 est noté $P$, et la base du théorème 2 de la page 22 est notée $B$ ;
les pages les appellent $X'$, $X'$ et $Y$.

La dimension d'une partie fermée de $X'$ est sa dimension comme espace
topologique. Pour $\mathfrak{p} \in X'$, l'adhérence de $\mathfrak{p}$ dans
$X'$ est de dimension $\dim A/\mathfrak{p} - 1$ : le spectre épointé décale
les dimensions d'un cran.

Un espace topologique noethérien $T$ est \emph{connexe en dimension $\geq k$}
si, pour toute partie fermée $Z \subset T$ de dimension $< k$, l'espace
$T \smallsetminus Z$ est connexe et non vide. Pour $k \leq 0$ cela veut dire
seulement que $T$ est connexe et non vide.\footnote{La notion est de
R.~Hartshorne, « Complete intersections and connectedness », \emph{Amer.\ J.\
Math.}\ 84 (1962) ; c'est sans doute lui que la page 6 cite. Les pages
emploient l'expression sans la définir, et les conditions $a_j$, $b_j$, $c_j$
de la page 8 ne sont définies nulle part dans le dossier. Nous fixons cette
convention, qui est celle qu'exigent les énoncés pour être vrais tels que nous
les donnons.}

Le corps de base est noté $\Bbbk$, la lettre $k$ restant l'entier de la
connexité.\footnote{Les pages 26 et 55 écrivent $k$ pour les deux.} Pour une
partie fermée $Y$ d'un schéma localement noethérien $X$ et un Module cohérent
$F$, $X_{/Y}$ ou $\widehat{X}$ est le complété formel de $X$ le long de $Y$, et
$F_{/Y}$ ou $\widehat{F}$ le complété de $F$. Le faisceau structural est
$\mathcal{O}_X$.\footnote{Il souligne souvent le $\mathcal{O}$ ; on ne garde
pas ce soulignement.}

\section*{I. La connexité en dimension $\geq k$}

\pagerange{2}{6}
\subsection*{Deux lemmes sur le spectre épointé (pages 2 à 6)}

Soient $A$ un anneau local noethérien et $f \in \mathfrak{m}$. On pose
$Y' = X' \cap V(f)$. Le premier énoncé est le cas le plus simple du
théorème à venir.

\emph{Lemme 2.} Supposons $A$ complet, les composantes irréductibles de $X'$
de dimension $\geq 2$, et $X'$ connexe en dimension $\geq 1$. Alors les
composantes de $Y'$ sont de dimension $\geq 1$, et $Y'$ est
connexe.\footnote{L'énoncé est un palimpseste de quatre lignes biffées et
récrites ; seule l'hypothèse de dimension se lit avec quelque assurance, et «
complet », encerclé, est ajouté en marge. L'hypothèse de connexité en
dimension $\geq 1$ est fournie par nous : la page 10 dit que le lemme est le
cas $k = 1$ du théorème de la page 8, qui la contient. Elle est nécessaire :
si $X$ est la réunion de deux $3$-plans de $\Bbbk^6$ qui se coupent en
l'origine, $X'$ a deux composantes disjointes de dimension $2$, et une section
générale $Y'$ les coupe en deux morceaux disjoints.}

La démonstration est presque entièrement illisible ; on y voit l'argument de
dimension sur lequel elle repose : pour une composante $Z'$ de $X'$, chaque
composante non vide de $Z' \cap V(f)$ est de dimension $\geq \dim Z' - 1$, et
$Z' \cap V(f)$ n'est pas vide dès que $\dim Z' > 0$ --- c'est le théorème de
l'idéal principal de Krull, appliqué dans $\overline{Z'}$ qui contient
$\mathfrak{m}$.\footnote{La seconde note marginale de la page 2 dit « si $Z'
\cap Y' \neq \emptyset$ si $\dim Z' > 0$ ». Un mot encerclé, lu «
fermée-ouverte », suggère que la suite raisonne sur une partie ouverte et
fermée de $Y'$.}

Un \emph{Corollaire 1}, aux pages 2 et 4, passe à la clôture normale
$\tilde{A}$ d'un $A$ intègre de dimension $\geq 2$ et compare les sections de
$\mathcal{O}$ sur $X'$ et sur un complété ; son état ne permet pas de le
restituer.\footnote{On lit seulement « $\tilde{A}$ clôture normale … fini … »,
« de dim $\geq 3$ » et la formule $H^0(\widehat{X}', \mathcal{O}) \simeq
H^0(X', \mathcal{O})$, dont les indices et l'accent sont lus d'après le
tracé.}

Le second lemme fournit, pour $f$ assez général, ce que le théorème donnera
pour tout $f$.

\emph{Lemme 3.} Soit $k \geq 1$. Supposons $A$ complet, les composantes de $X'$ de
dimension $\geq k+1$ et $X'$ connexe en dimension $\geq k$. Il existe une
partie finie $\Phi \subset X'$ telle que, pour tout $f \in \mathfrak{m}$
vérifiant $V(f) \cap \Phi = \emptyset$, le fermé $Y' = X' \cap V(f)$ soit
connexe en dimension $\geq k-1$.\footnote{La forme de l'énoncé vient de la
marge, qui relie « il existe $\Phi$ fini en $X'$ … $V(f) \cap \Phi =
\emptyset$ » au « Alors » ; le corps de l'énoncé est biffé et récrit.}

La démonstration (pages 4 et 6) se ramène, par un Lemme 1 et un Corollaire 2
qui ne sont pas dans le dossier, au cas où $X$ est irréductible et normal ;
elle choisit $\Phi$ fini dans $X'$, évité par $V(f)$, et conclut par la
condition $(S_2)$ de Serre et le théorème de Hartshorne --- un anneau local
noethérien qui vérifie $(S_2)$ a un spectre connexe en codimension
$1$.\footnote{Seuls se lisent les formules, « $(S_2)$ » et le renvoi, incertain,
à « Hartshorne ». Nous ne reconstituons pas l'argument. La page ne suppose
pas $A$ complet ; nous n'affirmons l'énoncé que dans ce cas, où il est de
toute façon contenu dans le théorème qui suit, qui vaut pour tout $f$ : le
lemme en est l'étape de récurrence.}

\pagerange{8}{12}
\subsection*{Le théorème de connexité local (pages 8 à 12)}

\emph{Théorème.} Soient $A$ un anneau local noethérien, $\widehat{A}$ son
complété, $\widehat{X}'$ le spectre épointé de $\widehat{A}$, et $k \geq 1$.
Supposons que
\begin{itemize}
\item[(i)] les composantes irréductibles de $\widehat{X}'$ soient de dimension
$\geq k+1$ ;
\item[(ii)] $\widehat{X}'$ soit connexe en dimension $\geq k$.
\end{itemize}
Alors, pour $0 \leq l \leq k$ et $f_1, \ldots, f_l \in \mathfrak{m}$, le fermé
\[
  Y' = X' \cap V(f_1) \cap \cdots \cap V(f_l)
\]
a ses composantes de dimension $\geq k+1-l$ et il est connexe en dimension
$\geq k-l$. En particulier, pour $l = k$, il est connexe et non vide.

C'est, sous sa forme locale, le \emph{théorème de connexité de Grothendieck},
tel qu'il figure à l'exposé XIII de SGA~2.\footnote{Le Théorème est encadré.
La page met l'hypothèse de dimension sur $\widehat{A}$ --- en notant entre
crochets qu'il suffit pour cela que les composantes de $X$ soient de dimension
$\geq k+2$, sous une condition sur $A \to \widehat{A}$ illisible --- mais
l'hypothèse de connexité sur $X'$. Nous mettons les deux sur le complété, ce
qui revient à supposer $A$ complet ; c'est la forme où l'énoncé est démontré,
et nous n'affirmons pas la variante de la page. Les conclusions, elles,
redescendent de $\widehat{A}$ à $A$ : $\operatorname{Spec}\widehat{A} \to
\operatorname{Spec} A$ est surjectif et conserve la dimension des fermés, et
les composantes de $\operatorname{Spec}(\widehat{A}/\mathfrak{p}\widehat{A})$,
pour $\mathfrak{p}$ premier minimal de $A$, sont des composantes de
$\operatorname{Spec}\widehat{A}$ (par platitude). Les conditions
« $a_{k-l}$, $b_{k-l}$, $c_{k-l}$ » de la page ne sont pas définies ; nous
lisons les deux premières comme les deux conclusions ci-dessus et ne savons
pas ce qu'est la troisième.} Les hypothèses et la borne sur $l$
sont nécessaires : la réunion de deux $3$-plans de $\Bbbk^6$ qui se coupent en
l'origine vérifie (i) pour $k = 1$ mais non (ii), et l'on a vu que sa section
par un hyperplan général n'est pas connexe ; un cône sur une conique vérifie (i) et (ii) pour
$k = 0$, et sa section par un hyperplan général passant par le sommet a deux
composantes.

\emph{Réduction au cas $l = k$} (pages 8 et 10). On peut supposer $A$ complet.
L'assertion sur les dimensions est le théorème de l'idéal principal dans
l'anneau caténaire $\widehat{A}$. Si $Y'$ était déconnecté par un fermé $Z'$
de dimension $< k-l$, on choisirait $f_{l+1}, \ldots, f_k \in \mathfrak{m}$
tels que $Z' \cap V(f_{l+1}) \cap \cdots \cap V(f_k) = \emptyset$ ; chacun des
deux morceaux de $Y' \smallsetminus Z'$ contient une partie de dimension
$\geq k+1-l$ et rencontre donc $V(f_{l+1}, \ldots, f_k)$, de sorte que
$X' \cap V(f_1, \ldots, f_k)$ serait non connexe. Il suffit donc de montrer :

\emph{Corollaire 1} (page 10). Sous les hypothèses du théorème, pour
$f_1, \ldots, f_k \in \mathfrak{m}$, le fermé
$X' \cap V(f_1) \cap \cdots \cap V(f_k)$ est connexe et non
vide.\footnote{L'énoncé, souligné, n'est lisible que dans ses formules ; la
conclusion « connexe » est celle qu'appelle la démonstration qui suit, dont le
mot clé, deux fois, est « discontinu ».}

\emph{Démonstration} (pages 10 et 12), par récurrence sur $k$. Pour $k = 1$,
c'est le Lemme 2. Soit $k \geq 2$, l'énoncé connu jusqu'à $k-1$. Montrons
d'abord que $X' \cap V(f_1)$ est connexe en dimension $\geq k-1$. Sinon, soit
$Z'$ un fermé de dimension $< k-1$ qui le déconnecte. Soit $\Phi$ la partie
finie que donne le Lemme 3. Par évitement des idéaux premiers, on choisit
$f_2, \ldots, f_k \in \mathfrak{m}$ tels que $Z' \cap V(f_1, \ldots, f_k) =
\emptyset$ et $V(f_2) \cap \Phi = \emptyset$. Le Lemme 3 dit que
$X' \cap V(f_2)$ vérifie les hypothèses au rang $k-1$ ; l'hypothèse de
récurrence, appliquée aux $k-1$ éléments $f_1, f_3, \ldots, f_k$, dit que
$X' \cap V(f_1, \ldots, f_k)$ est connexe --- alors que le choix de $f_2,
\ldots, f_k$ le rend non connexe, comme dans la réduction. Ainsi
$X' \cap V(f_1)$ vérifie les hypothèses au rang $k-1$, et la récurrence,
appliquée à $f_2, \ldots, f_k$, achève.\footnote{La page 12 porte « Or »
souligné trois fois devant l'appel au Lemme 3. Le texte de liaison est presque
entièrement illisible ; la démonstration est restituée d'après ses formules,
et son ordre est celui que celles-ci imposent.}

\pagerange{12}{16}
\subsection*{Une réciproque (pages 12 à 16)}

Un \emph{Corollaire I} énonce une réciproque : si, pour tous
$f_1, \ldots, f_k \in \mathfrak{m}$, l'ensemble
$\pi_0\bigl(X' \cap V(f_1) \cap \cdots \cap V(f_k)\bigr)$ des composantes
connexes est réduit à un point, alors les composantes de $X'$ sont de
dimension $\geq k+1$ et $X'$ est connexe en dimension $\geq k$. Les pages 14 et
16, qui le démontrent, sont en grande partie annulées par de longs traits et
des cadres barrés ; on y voit le principe --- si une composante de $X'$ est de
dimension $l \leq k$, des équations $f_1, \ldots, f_l$ convenables la font
disparaître ou la séparent du reste --- et un « !! » final. Nous ne
reconstituons pas l'énoncé exact et ne l'affirmons pas.\footnote{L'hypothèse
du Corollaire I est presque illisible ; « $(k \geq 1)$ » est ajouté au-dessus
du titre, et une note en marge, en partie biffée, demande d'« exclure » le cas
de $X$ irréductible de dimension $1$ --- celui où $X' \cap V(f)$ est vide pour
tout $f$ non diviseur de zéro, et où la question de savoir si le vide est
connexe se pose. Un exemple montre comment la réciproque opère : pour $X$
réunion d'un $3$-plan et d'un plan qui se coupent selon une droite, et $k = 1$,
la composante issue du plan est de dimension $1$ dans $X'$, et pour $f$ général
$X' \cap V(f)$ est la réunion disjointe d'une surface et d'une droite
épointées, qui ne se rencontrent qu'en $\mathfrak{m}$.}

\pagerange{18}{18}
\subsection*{Le groupe fondamental (page 18)}

\emph{Corollaire.} Soit $A$ un anneau local noethérien complet, normal, de
dimension $\geq n+1$ (donc $\dim X' \geq n$), et soient
$f_1, \ldots, f_k \in \mathfrak{m}$ avec $k \leq n-1$. Alors
$Y' = X' \cap V(f_1) \cap \cdots \cap V(f_k)$ est connexe et non vide, et
l'homomorphisme $\pi_1(Y') \to \pi_1(X')$ des groupes fondamentaux étales
(pour un point de base dans $Y'$) est surjectif.\footnote{L'hypothèse sur
$X'$, devant « de dim $\geq n$ », est illisible ; « surjectif » est la
lecture de la conclusion. Nous fournissons « complet et normal » : la
démonstration utilise que les revêtements connexes de $X'$ sont irréductibles
(« Or $X'_1$ est irréductible »), ce que donne la normalité, et le théorème
précédent demande la complétude. Les « $H_i$ » de la page sont les $V(f_i)$.}

En effet, $\pi_1(Y') \to \pi_1(X')$ est surjectif si et seulement si tout
revêtement étale connexe $X'_1 \to X'$ reste connexe au-dessus de $Y'$. Un tel
$X'_1$ est normal et connexe, donc intègre ; la normalisation $B$ de $A$ dans
son corps des fonctions est finie sur $A$ (qui est complet), locale (car $A$
est hensélien et $B$ intègre), normale, de même dimension, et $X'_1$ est son
spectre épointé. Le théorème, appliqué à $B$ avec $k$ remplacé par $n-1$,
donne la connexité de $X'_1 \cap V(f_1, \ldots, f_k)$ ; appliqué à $A$, celle
de $Y'$. C'est une forme du \emph{théorème de Lefschetz local pour le
$\pi_1$}.\footnote{Suit un paragraphe « Conséquences », sur le $\pi_1$ du
spectre épointé d'un anneau local et le cas de la dimension $2$, presque
entièrement illisible et terminé par un point d'exclamation.}

\pagerange{20}{20}
\subsection*{Du local au global (page 20)}

\emph{Lemme 4.} Soient $X$ un schéma localement noethérien, connexe, dont les
composantes irréductibles sont de dimension $\geq k$, et supposons que pour
tout point $z \in X$ tel que $\dim \overline{\lbrace z \rbrace} < k$, le
spectre épointé de $\mathcal{O}_{X,z}$ soit connexe en dimension
$\geq k - 1 - \dim \overline{\lbrace z \rbrace}$. Alors $X$ est connexe en
dimension $\geq k$.\footnote{La page demande que chaque
$(\operatorname{Spec} \mathcal{O}_{X,x})'$ soit connexe en dimension $\geq k$ :
c'est plus fort que l'hypothèse ci-dessus, et le lemme vaut donc aussi sous la
forme de la page. Nous fournissons la connexité de $X$, sans laquelle
l'énoncé est faux (une réunion disjointe) : l'adjectif qui suit « préschéma »
est illisible. La note marginale, presque illisible, renvoie à EGA~IV.}

\emph{Démonstration.} Soit $Z$ un fermé de dimension $< k$ tel que
$X \smallsetminus Z = U' \sqcup U''$, avec $U'$, $U''$ ouverts, fermés et non
vides. Aucune composante de $X$ n'est contenue dans $Z$, donc
$X = \overline{U'} \cup \overline{U''}$ ; comme $X$ est connexe, il existe
$z \in \overline{U'} \cap \overline{U''}$, et nécessairement $z \in Z$. Dans
$X_z = \operatorname{Spec} \mathcal{O}_{X,z}$, le spectre épointé privé de
$Z_z$ est la réunion disjointe des traces de $U'$ et de $U''$, toutes deux non
vides puisque $z$ est adhérent à chacun. Or les points de $Z_z$ autres que $z$
ont dans $X_z$ des adhérences de dimension
$\leq \dim \mathcal{O}_{Z,z} - 1 \leq \dim Z - \dim \overline{\lbrace z
\rbrace} - 1 < k - 1 - \dim \overline{\lbrace z \rbrace}$, ce qui contredit
l'hypothèse.\footnote{La page localise de même en un point $z$ de
$\overline{U'} \cap \overline{U''}$ et conclut « absurde ». Dans « $X^{z} =
X' \cup X''$ », l'exposant est douteux ; nous lisons $X = \overline{U'} \cup
\overline{U''}$.}

\pagerange{22}{26}
\subsection*{Les fibres spéciales d'un morphisme propre (pages 22 à 26)}

\emph{Théorème 2.} Soient $f\colon X \to B$ un morphisme propre de schémas
localement noethériens, avec $B$ intègre et excellent, $b_0 \in B$, $\eta$
le point générique de $B$, $X_0$ et $X_1$ les fibres en $b_0$ et en $\eta$.
Supposons :
\begin{itemize}
\item[(a)] $B$ géométriquement unibranche en $b_0$, et les composantes
irréductibles de $X$ dominantes sur $B$ ;
\item[(b)] les composantes irréductibles de $X_1$ de dimension $\geq k+1$, et
$X_1$ \emph{géométriquement} connexe en dimension $\geq k$.
\end{itemize}
Alors les composantes irréductibles de $X_0$ sont de dimension $\geq k+1$, et
$X_0$ est connexe en dimension $\geq k$.\footnote{Trois écarts. (1) La page
écrit « de dim $\geq k$ » dans (b), lecture serrée en bout de ligne ; la
conclusion, la page 24 (« $X_1$ irréductible de dim $\geq k+1$ ») et la formule
de la page 26 demandent $k+1$. (2) Le « géométriquement » de (b) est de nous.
Il est nécessaire déjà pour $k = 0$ : si $B$ est le spectre d'un anneau de
valuation discrète et $X$ celui du normalisé de $B$ dans un corps de nombres où
l'idéal premier se décompose, $X_1$ est un point, connexe, et $X_0$ en a deux.
Celui de (a), entre crochets sur la page, est d'une lecture incertaine. (3)
L'excellence de $B$ est de nous : elle sert à passer, dans la démonstration,
à des anneaux locaux dont le complété reste intègre.}

L'assertion sur les dimensions est la semi-continuité de la dimension des
fibres ; elle n'utilise pas l'hypothèse unibranche. Celle-ci sert, comme le
dit la page, au \emph{théorème de connexion} de Zariski, qui assure que $X_0$
est connexe.\footnote{Pour $k = 0$ l'énoncé est exactement ce théorème (EGA~III,
§~4.3, et sa variante à base unibranche dans EGA~IV).} Pour la connexité en
dimension $\geq k$, le Lemme 4 ramène à l'énoncé local :

\emph{Lemme 5.} Sous les hypothèses du théorème 2, pour $x \in X_0$, le spectre
épointé de $\mathcal{O}_{X_0,x}$ est connexe en dimension
$\geq k - 1 - \dim \overline{\lbrace x \rbrace}$ (adhérence prise dans
$X_0$).\footnote{La page énonce la connexité « en dim $\geq k$ », avec une note
marginale où paraît « $\dim \overline{x}$ » ; nous donnons la borne que le
Lemme 4, tel que nous l'énonçons, consomme.}

L'esquisse des pages 24 et 26 est la suivante. Par le Lemme 1 (absent du
dossier), on se ramène à $X_1$ irréductible, puis $X$ intègre. Pour $x$ fermé
dans $X_0$, la formule des dimensions donne
\[
  \dim \mathcal{O}_{X,x} = \dim \mathcal{O}_{B,b_0} + \dim X_1 \geq n + k + 1,
  \qquad n = \dim \mathcal{O}_{B,b_0} ;
\]
la fibre $X_0$ est, au voisinage de $x$ et ensemblistement, définie par les $n$
éléments d'un système de paramètres de $\mathcal{O}_{B,b_0}$, et le théorème
local, appliqué avec $K = n + k - 1$ à $\mathcal{O}_{X,x}$ et à ces $n$
équations, donne la connexité en dimension $\geq K - n = k - 1$, ce que le
Lemme 5 demande en un point fermé. Nous ne donnons ce raisonnement que comme une esquisse : il demande que
le spectre épointé du complété de $\mathcal{O}_{X,x}$ soit connexe en
dimension $\geq \dim \mathcal{O}_{X,x} - 2$, ce qui vaut par exemple si ce
complété est intègre --- d'où le passage au normalisé de $X$, permis par
l'excellence --- et ce que la page tient pour « facile ».\footnote{La page
écrit $\dim \mathcal{O}_{X,x} = \dim \mathcal{O}_{Y,y_0} + \dim X_1$ et
« $\geq K+1 = (n+k)+1$ » puis « connexe en dim $\geq K - n = k$ », $K$
étant une capitale nette : elle compte les dimensions dans
$\operatorname{Spec} \mathcal{O}_{X,x}$, nous dans son spectre épointé, d'où
le décalage d'un cran. Nous n'avons pas
vérifié l'esquisse dans tous ses détails ; l'énoncé du théorème~2 est donné
sous les hypothèses où elle nous paraît aller jusqu'au bout.}

Trois corollaires suivent. Les deux premiers tiennent en une ligne : une
variante pour les fibres géométriques, et une variante pour $f$ propre et
ouvert. Le troisième est le théorème de connexité sous sa forme projective :

\emph{Corollaire 3.} Soit $X \subset \mathbf{P}^N_{\Bbbk}$ un fermé dont les
composantes sont de dimension $\geq k+1$ et qui est connexe en dimension
$\geq k$. Pour des hyperplans $H_1, \ldots, H_\ell$ avec $\ell \leq k$, le
fermé $X \cap H_1 \cap \cdots \cap H_\ell$ est connexe en dimension
$\geq k - \ell$.

On le déduit du théorème local appliqué au cône affine de $X$ en son sommet :
l'anneau local est gradué, son complété a pour composantes les complétés des
cônes sur les composantes de $X$, et le spectre épointé décale les dimensions
exactement comme le passage de $X$ à son cône les relève. C'est l'énoncé que
la littérature attribue à Grothendieck sous ce nom.\footnote{Voir W.~Fulton et
R.~Lazarsfeld, « Connectivity and its applications in algebraic geometry »,
Lecture Notes in Math.\ 862 (1981), qui le citent d'après SGA~2, XIII. Une
note marginale de la page 26 annonce une « généralisation au cas … d'un
morphisme $X \to \mathbf{P}^r_{\Bbbk}$ » ; c'est la forme que W.~Fulton et
J.~Hansen (1979) et P.~Deligne ont donnée au théorème de connexité. La marge
est trop illisible pour dire ce qu'il visait, et rien n'indique qu'il ait eu
leur énoncé.}

\section*{II. Les couples de Lefschetz-Grauert}

\pagerange{29}{31}
\subsection*{Introduction et plan (pages 29 à 31)}

La page 29 porte « III » et « Introduction » : c'est le début d'un chapitre
dont le dossier ne contient pas ce qui précède. Elle rappelle qu'au
paragraphe précédent, pour $X$ propre sur un schéma localement noethérien $Y$
et $Y'$ fermé dans $Y$, on a comparé les foncteurs $R^i f$ et $R^i \hat{f}$
--- c'est le théorème des fonctions formelles --- et que, sur la base d'un
anneau complet, un théorème d'existence donne l'équivalence entre Modules
cohérents sur $X$ et Modules cohérents sur le complété formel : c'est le
contenu des paragraphes 4 et 5 d'EGA~III. Le chapitre annonce des
généralisations à des morphismes non propres --- une immersion ouverte, par
exemple --- au moyen d'un théorème de finitude et de la notion de
\emph{profondeur}, et en fait l'outil d'une étude systématique des propriétés
« du type de Lefschetz » : relier les propriétés de $X$ à celles d'un
sous-schéma $Y$ \emph{via} le complété formel $X_{/Y}$. Il cite H.~Grauert
(« 1958 ») et des théorèmes de Serre.\footnote{Les pages 29 et 30 sont un
palimpseste rapide ; on n'en restitue que l'articulation. Les références sont
des crochets vides, et celle à Grauert ne s'identifie pas d'après la page ;
nous n'en proposons pas. Le nom « Lefschetz-Grauert » n'est pas celui que
SGA~2 retiendra.}

Parmi les résultats annoncés : pour $X$ projectif, $Y$ section hyperplane et
$F$ cohérent dont la profondeur aux points fermés de $X \smallsetminus Y$ est
$\geq p$, l'homomorphisme
\[
  H^i(X, F) \longrightarrow H^i(X_{/Y}, F_{/Y})
\]
est bijectif pour $i \leq p - 2$ (et injectif pour $i = p-1$).\footnote{La page
a « bijectif [… $i \leq \operatorname{prof} F - 2$] » pour $F$ localement
libre ; le cadre --- section hyperplane projective --- est lu dans la phrase
précédente. Nous donnons l'énoncé sous la forme de SGA~2, XII. Un second
point (ii), où paraît « $\operatorname{prof}(\mathcal{O}_{\widehat{X}}) \geq
2$ », est illisible.}

Le plan de la page 31, en neuf points : 1.~généralisation des théorèmes de
finitude formelle et de comparaison, application à un théorème d'existence ;
2.~le théorème de finitude locale ; 3.~cas particulier des anneaux locaux ;
4.~cas particulier des schémas projectifs ; 5.~fermés de Lefschetz-Grauert ;
6.~propriétés liées à $\pi_0$, $\pi_1$, $\operatorname{Pic}$ ;
7.~permanence ; 8.~théorèmes de pureté et théorèmes de Lefschetz pour le
$\pi_1$ ; 9.~anneaux quasi-factoriels et théorème de Lefschetz pour
$\operatorname{Pic}$.\footnote{Deux anciens points 3 et 4 sont barrés et
renumérotés. Une note encadrée relie les points 6 et 8 : « utiliser IV, § …
sur les morphismes étales », sans doute EGA~IV.} Le dossier développe les
points 5 à 7 (pages 32 à 51), un morceau du point 4 (pages 55 à 61) et
quelques calculs pour le point 9 (pages 64 à 66).

\pagerange{32}{36}
\subsection*{Définitions et premières remarques (pages 32 à 36)}

Soient $X$ un schéma localement noethérien, $Y$ une partie fermée, $n$ un
entier. Tous les Modules localement libres sont de type fini.

\emph{Définition.} On dit que $(X, Y)$ vérifie $(\mathrm{LG})_n$ si, pour
tout Module localement libre $F$ sur $X$,
\[
  (*) \qquad H^i(X, F) \longrightarrow H^i(X_{/Y}, F_{/Y})
\]
est bijectif pour $i \leq n$ et injectif pour $i = n+1$. On dit que $(X, Y)$
vérifie \emph{LG faible} si $(*)$ est bijectif pour $i = 0$ et tout $F$
localement libre, c'est-à-dire si $F \mapsto F_{/Y}$ est pleinement fidèle des
Modules localement libres sur $X$ dans ceux sur $X_{/Y}$. On dit que $(X, Y)$
vérifie \emph{LG} si $(U, Y)$ vérifie LG faible pour tout voisinage ouvert $U$
de $Y$, et \emph{LG stricte} si de plus tout Module localement libre sur
$X_{/Y}$ est le complété d'un Module localement libre sur un voisinage de $Y$.
Notant $\mathcal{L}(U)$ la catégorie des Modules localement libres sur $U$,
LG stricte signifie que les foncteurs de restriction $\mathcal{L}(U) \to
\mathcal{L}(V)$, pour $U \supset V \supset Y$, sont pleinement fidèles et que
\[
  \varinjlim_{U \supset Y} \mathcal{L}(U) \longrightarrow \mathcal{L}(X_{/Y})
\]
est une équivalence de catégories.\footnote{Il écrit sous la limite
« "Pseudolimite" de catégories » : c'est la $2$-colimite. Les mots « faible »
et « stricte » sont d'une lecture incertaine ; la page 35 dit aussi «
effective ». La définition de LG est restituée d'après la Remarque 3 et la
conclusion de la page 36, la page 32 étant illisible à cet endroit.} Ce sont,
sous les noms de condition de Lefschetz $\mathrm{Lef}(X, Y)$ et de condition
de Lefschetz effective $\mathrm{Leff}(X, Y)$, les notions de l'exposé X de
SGA~2.\footnote{« et injectif pour $i = n+1$ » est encerclé dans la
définition, et la marge porte « inutile » : il s'est demandé si cette clause
était superflue ; c'est la Question de la Remarque 1. Deux notes encadrées
de la marge posent l'une une conséquence (« si $\overline{\mathfrak{z}} \cap
Y = \emptyset$, alors $\mathfrak{z} \notin \operatorname{Ass}
\mathcal{O}_X$ », que démontre le Lemme de la page 36), l'autre une question
restée ouverte sur la page : LG implique-t-il $\operatorname{prof}
\mathcal{O}_{X,x} \geq 2$ pour $\overline{\lbrace x \rbrace} \cap Y =
\emptyset$ ?}

\emph{Remarque 1} (page 33). LG faible équivaut à $(\mathrm{LG})_0$ : la
bijectivité en degré $0$ entraîne l'injectivité en degré $1$. En effet, une
classe de $H^1(X, F) = \operatorname{Ext}^1(\mathcal{O}_X, F)$ est celle d'une
extension $0 \to F \to E \to \mathcal{O}_X \to 0$, avec $E$ localement libre ;
si elle s'annule sur $X_{/Y}$, la section formelle de $\widehat{E}$ qui scinde
l'extension complétée provient d'une section $s$ de $E$, dont l'image dans
$H^0(X, \mathcal{O}_X)$ a pour complété $1$, donc vaut $1$ : l'extension est
scindée.

La page pose ensuite la question générale --- si $(*)$ est bijectif pour
$i \leq n$ et tout $F$, est-il injectif pour $i = n+1$ ? --- et répond « Oui
», en plongeant $F$ dans un autre Module. Nous savons le démontrer lorsque $X$
est projectif sur un anneau noethérien, et $n \geq 0$ : on plonge $F$ dans
$L = \mathcal{O}_X(m)^{\oplus N}$, à conoyau $G$ localement libre, avec
$m$ assez grand pour que $H^{n+1}(X, L) = 0$. Une classe $c \in H^{n+1}(X, F)$
de complété nul a une image nulle dans $H^{n+1}(X, L)$, provient donc de
$g \in H^n(X, G)$ ; le complété de $g$ provient de $H^n(\widehat{L}) =
H^n(L)^{\wedge}$, et l'injectivité en degré $n$ pour $G$ montre que $g$ provient
de $H^n(X, L)$, d'où $c = 0$.\footnote{Après « Oui, prendre une injection
$0 \to F \to I \to G \to 0$ », la page est en interligne serré et illisible.
La démonstration ci-dessus est la nôtre, et l'hypothèse projective aussi ;
nous ne savons pas si la réponse est « oui » en général.}

\emph{Remarque 2} (pages 33 et 34). Soit $(L_\alpha)$ une famille de Modules
localement libres telle que tout Module localement libre $F$ se plonge dans
une somme directe finie $L$ de $L_\alpha$, avec un conoyau localement libre
--- ce qui revient à dire que $F^\vee$ est quotient d'une somme de
$L_\alpha^\vee$. Exemples : $X$ quasi-affine et la famille réduite à
$\mathcal{O}_X$ ; $X$ quasi-projectif sur un anneau noethérien et les
$\mathcal{O}_X(m)$, $m \geq m_0$. Alors $(\mathrm{LG})_n$ équivaut à : $(*)$
est bijectif pour $i \leq n$ et injectif pour $i = n+1$ lorsque $F$ est l'un
des $L_\alpha$.

Par récurrence sur $n$, le cas $n = -1$ étant clair. Supposons
$(\mathrm{LG})_{n-1}$ et soit $0 \to F \to L \to G \to 0$ comme ci-dessus. On
compare les suites de cohomologie :

\begin{tikzcd}[column sep=small, row sep=small, nodes={font=\scriptsize}]
H^i(L) \arrow[r] \arrow[d, "\alpha"'] & H^i(G) \arrow[r] \arrow[d, "\beta"] & H^{i+1}(F) \arrow[r] \arrow[d, "\varphi"] & H^{i+1}(L) \arrow[r] \arrow[d, "\alpha'"] & H^{i+1}(G) \arrow[d, "\beta'"] \\
H^i(\widehat{L}) \arrow[r] & H^i(\widehat{G}) \arrow[r] & H^{i+1}(\widehat{F}) \arrow[r] & H^{i+1}(\widehat{L}) \arrow[r] & H^{i+1}(\widehat{G})
\end{tikzcd}

Pour $i = n-1$ : $\beta$ est bijectif et $\beta'$ injectif (récurrence),
$\alpha$ et $\alpha'$ bijectifs (hypothèse) ; le lemme des cinq donne $\varphi$
bijectif, c'est-à-dire $(*)$ bijectif en degré $n$ pour tout $F$. Pour $i = n$ :
$\alpha$ est surjectif, $\beta$ injectif (ce qu'on vient de voir, appliqué à
$G$), $\alpha'$ injectif ; donc $\varphi$ est injectif.\footnote{La page
conclut « $\varphi$ est surjection (bijection) » au cas $\alpha$) et esquisse
le cas $\beta$) ; ses primes ne se distinguent pas, et son dernier symbole est
lu $\varphi'$ avec doute. Le crochet « [N.B. … !] » est biffé.}

\emph{Remarque 3} (pages 35 et 36). La condition LG passe d'elle-même à tout
voisinage ouvert de $Y$ ; il n'en va pas de même en général de
$(\mathrm{LG})_0$, car un Module localement libre sur un voisinage $U$ de $Y$
ne s'étend pas toujours à $X$. Mais si tout Module localement libre sur $U$
est sous-module, à conoyau localement libre, de la restriction d'un Module
localement libre sur $X$ --- c'est le cas si $X$ est quasi-projectif ---,
alors $(\mathrm{LG})_0$ pour $(X, Y)$ entraîne $(\mathrm{LG})_0$ pour $(U, Y)$,
c'est-à-dire LG. La surjectivité de $H^0(U, F) \to H^0(X_{/Y}, F_{/Y})$ se
déduit de celle sur $X$ par le même plongement ; l'injectivité tient à ce
qu'une section de complété nul est nulle au voisinage de $Y$ (théorème
d'intersection de Krull), donc à support dans un fermé de $U$ disjoint de
$Y$, et au lemme suivant.

\emph{Lemme} (page 36). Si $X$ est quasi-projectif sur un anneau noethérien
et si $(X, Y)$ vérifie $(\mathrm{LG})_0$, alors tout point $x \in X$ tel que
$\overline{\lbrace x \rbrace} \cap Y = \emptyset$ n'est pas associé à
$\mathcal{O}_X$.

Sinon, $\mathcal{O}_X$ contiendrait un sous-faisceau cohérent non nul $F$ à
support dans $\overline{\lbrace x \rbrace}$ ; pour $m$ assez grand, $F(m)
\subset \mathcal{O}_X(m)$ aurait une section non nulle, dont le complété le
long de $Y$ est nul : cela contredit l'injectivité de $(*)$ pour
$\mathcal{O}_X(m)$.

\emph{Conclusion pratique.} Si $X$ est quasi-projectif sur un anneau
noethérien,
\[
  (\mathrm{LG})_0 \iff \text{LG faible} \iff \text{LG} \iff
  H^0(X, \mathcal{O}_X(m)) \xrightarrow{\ \sim\ } H^0(X_{/Y}, \mathcal{O}_X(m)_{/Y})
  \ \text{pour } m \gg 0 .
\]
La dernière condition entraîne LG faible par l'argument de la Remarque 2 en
degré $0$, qui n'utilise que la bijectivité de $H^0$ sur les
$\mathcal{O}_X(m)$.\footnote{La conclusion est encadrée. La page écrit une
implication vers le quatrième membre, et trois flèches courbes au crayon,
marquées « $n$ », remontent vers les précédents ; l'équivalence est vérifiée
ici. Le troisième membre se lit « $(\mathrm{LG})_X$ », indice incertain ; nous
lisons LG, ce que la Remarque 3 établit.}

\pagerange{37}{40}
\subsection*{Composantes connexes, groupe fondamental, groupe de Picard (pages 37 à 40)}

Le complété $X_{/Y}$ a le même espace sous-jacent que $Y$ ; ses idempotents
sont ceux de $\mathcal{O}_Y$, et ses revêtements étales sont ceux de $Y$
(SGA~1). Tout ce qui suit consiste à faire passer ces informations de
$X_{/Y}$ à $X$ par les conditions précédentes. On suppose $X$ noethérien.

\emph{Proposition} (page 37). Si $H^0(X, \mathcal{O}_X) \to H^0(X_{/Y},
\mathcal{O}_{X/Y})$ est injectif (resp.\ bijectif), alors $\pi_0(Y) \to
\pi_0(X)$ est surjectif (resp.\ bijectif). En particulier
$(\mathrm{LG})_{-1}$ (resp.\ $(\mathrm{LG})_0$) l'entraîne.

Une composante connexe de $X$ qui ne rencontre pas $Y$ a un idempotent non nul
de complété nul ; et si $H^0$ est bijectif, toute partition de $Y$ en ouverts
fermés, qui est un idempotent de $H^0(X_{/Y}, \mathcal{O})$, provient d'un
idempotent de $H^0(X, \mathcal{O}_X)$.

\emph{Lemme} (page 37). Supposons LG, et soit $F$ un Module localement libre
sur $X$. Toute structure d'algèbre sur $\widehat{F}$ provient d'une unique
structure d'algèbre sur $F$, et les décompositions de $\widehat{F}$ en
produit correspondent à celles de $F$ ; si $X' = \operatorname{Spec}(F)$,
$X'$ est connexe si et seulement si son complété l'est.\footnote{Le paragraphe
est barré de quatre longs traits, et « Lemme » remplace « Corollaire » biffé ;
il reste lisible dans ses formules, et la suite s'en sert. Une multiplication
est une section du Module localement libre
$\underline{\mathrm{Hom}}(F \otimes F, F)$, et les axiomes se vérifient après
complétion par injectivité de $(*)$.}

\emph{Proposition} (page 37). Supposons $X$ connexe, et prenons le point
de base des groupes fondamentaux dans $Y$. Si $(X, Y)$ vérifie LG,
$\pi_1(Y) \to \pi_1(X)$ est surjectif. Si $(X, Y)$ vérifie LG stricte,
l'homomorphisme
\[
  \pi_1(Y) \longrightarrow \varprojlim_{U} \pi_1(U),
\]
la limite étant prise sur les voisinages ouverts connexes $U$ de $Y$, est
bijectif.

\emph{Corollaire} (page 38). Si LG, le foncteur $X' \mapsto X' \times_X Y$
des revêtements étales de $X$ dans ceux de $Y$ est pleinement fidèle, et il
en va de même en remplaçant $X$ par un voisinage ouvert de $Y$. Si LG stricte,
pour tout revêtement étale $Y'$ de $Y$, il existe un voisinage ouvert $U$ de
$Y$ et un revêtement étale $U'$ de $U$ tels que $U' \times_U Y \simeq Y'$.

C'est le lemme appliqué aux algèbres localement libres qui définissent les
revêtements étales, et l'équivalence entre revêtements étales de $X_{/Y}$ et
de $Y$. La propriété sur les $\pi_1$ en est la traduction galoisienne. La page
annonce enfin des conditions « liées aux théorèmes de pureté » qui donnent
$\pi_1(Y) \simeq \pi_1(X)$, « Kif-Kif pour Pic » : c'est le point 8 du plan,
que le dossier ne développe pas.

\emph{Proposition} (page 39). (i) Si LG, $\operatorname{Pic}(X) \to
\operatorname{Pic}(X_{/Y})$ est injectif, et plus généralement les
$\operatorname{Pic}(U) \to \operatorname{Pic}(X_{/Y})$ le sont pour les
voisinages ouverts $U$ de $Y$, de sorte que les applications de transition
sont injectives et que
\[
  \varinjlim_{U \supset Y} \operatorname{Pic}(U) \longrightarrow \operatorname{Pic}(X_{/Y})
\]
est injectif. (ii) Si LG stricte, ce dernier homomorphisme est bijectif.

Si $\widehat{L}$ est trivial, une section formelle partout génératrice de
$\widehat{L}$ et son inverse proviennent de sections de $L$ et de
$L^{-1}$ dont le produit, de complété $1$, vaut $1$.

Supposons maintenant que $Y$ soit un diviseur de Cartier, d'Idéal
$\mathcal{J}$, et notons $\Lambda = \mathcal{O}_X(Y)$, de sorte que
$i^*(\Lambda)^{-1} \simeq \mathcal{J}/\mathcal{J}^2$, où $i\colon Y \to X$ est
l'injection. On écrit $\Lambda^{-n}$ pour $i^*(\Lambda)^{-n} \simeq
\mathcal{J}^n/\mathcal{J}^{n+1}$.

\emph{Lemme} (page 40). (i) Si $H^1(Y, \Lambda^{-n}) = 0$ pour $n \geq 1$,
l'homomorphisme $\operatorname{Pic}(X_{/Y}) \to \operatorname{Pic}(Y)$ est
injectif. (ii) Si de plus $H^2(Y, \Lambda^{-n}) = 0$ pour $n \geq 1$, il est
bijectif.

Pour $X_n$ le sous-schéma défini par $\mathcal{J}^{n+1}$, la suite exacte
$0 \to \mathcal{J}^n/\mathcal{J}^{n+1} \to \mathcal{O}_{X_n}^* \to
\mathcal{O}_{X_{n-1}}^* \to 1$ ($n \geq 1$) place dans
$H^1(Y, \Lambda^{-n})$ le noyau de $\operatorname{Pic}(X_n) \to
\operatorname{Pic}(X_{n-1})$ et l'obstruction à relever les sections
inversibles, et dans $H^2(Y, \Lambda^{-n})$ l'obstruction à relever un Module
inversible.\footnote{La page dit que, sous (i), le foncteur $\mathcal{F}
\mapsto \mathcal{F}_0$ des Modules inversibles sur $X_{/Y}$ dans ceux sur $Y$
est « pleinement fidèle », et une « équivalence » sous (ii). C'est trop :
deux automorphismes d'un Module inversible sur $X_{/Y}$ qui coïncident sur $Y$
diffèrent d'une unité $\equiv 1 \bmod \mathcal{J}$, et la fidélité demanderait
en outre $H^0(Y, \Lambda^{-n}) = 0$. Ce qui vaut tel quel est l'énoncé sur
$\operatorname{Pic}$, que la page en tire.}

\emph{Corollaire} (page 39). (i) Si $H^1(Y, \Lambda^{-n}) = 0$ pour $n \geq 1$
et si LG, alors $\operatorname{Pic}(X) \to \operatorname{Pic}(Y)$ est injectif,
ainsi que les $\operatorname{Pic}(U) \to \operatorname{Pic}(Y)$. (ii) Si
$H^1(Y, \Lambda^{-n}) = H^2(Y, \Lambda^{-n}) = 0$ pour $n \geq 1$ et si LG
stricte, alors $\varinjlim_{U \supset Y} \operatorname{Pic}(U)
\xrightarrow{\ \sim\ } \operatorname{Pic}(Y)$.\footnote{Au (ii), un trait
traverse « $H^1(Y, \Lambda^{-n}) = 0$ » ; ce n'est sans doute pas une biffure,
la conclusion demandant les deux annulations. Le corollaire est écrit avant le
lemme dont il découle.}

C'est le mécanisme du théorème de Lefschetz de Grothendieck pour le groupe de
Picard : sur une variété projective lisse de dimension $\geq 4$, la
restriction à une section hyperplane est un isomorphisme sur
$\operatorname{Pic}$.

\pagerange{40}{43}
\subsection*{Permanence (pages 40 à 43)}

Un \emph{Théorème 1} de la page 40, annulé, énonce la permanence de
$(\mathrm{LG})_n$ par un morphisme projectif et plat, puis un \emph{Th.~2}
s'interrompt au bas de la page ; la page 41 récrit le tout.

\emph{Théorème.} Soient $X$ un schéma noethérien quasi-projectif sur un anneau
noethérien, $Y$ une partie fermée, $f\colon W \to X$ un morphisme projectif,
$Y_W = f^{-1}(Y)$. Soit $F$ un Module cohérent sur $W$, à support propre sur
$X$ et \emph{plat} sur $X$, et considérons les homomorphismes
\[
  H^i(W, F) \longrightarrow H^i(W_{/Y_W}, F_{/Y_W}) .
\]
(i) Si $(X, Y)$ vérifie $(\mathrm{LG})_n$, ils sont bijectifs pour $i \leq n$
et injectifs pour $i = n+1$.

(ii) Si $(X, Y)$ vérifie LG, alors pour tout voisinage ouvert $U$ de $Y$, le
foncteur $F \mapsto F_{/Y_W}$, des Modules cohérents sur $f^{-1}(U)$ plats sur
$U$ à support propre dans les Modules cohérents sur $W_{/Y_W}$, est pleinement
fidèle ; plus précisément
$\operatorname{Hom}(F, G) \simeq \operatorname{Hom}(\widehat{F}, \widehat{G})$
dès que $G$ est plat sur $X$.

(iii) Si $(X, Y)$ vérifie LG stricte, tout Module cohérent sur $W_{/Y_W}$
plat sur $X_{/Y}$ est, sur un voisinage $f^{-1}(U)$ de $Y_W$, le complété
d'un Module cohérent plat sur $U$.\footnote{« plat sur $Y$ », dit la page à
l'hypothèse ; la lettre se lit nettement, mais (i) n'a de sens qu'avec « plat
sur $X$ », qui est ce que demandent (ii) et (iii). « quasi-projectif »,
doublement souligné, et « à support propre sur $X$ » sont des ajouts.
La précision de (ii) sur $\operatorname{Hom}(F, G)$ est en marge.}

Le manuscrit ne donne pas de démonstration. Pour (i), elle va ainsi : $F$
étant plat et à support propre, $Rf_*F$ est un complexe parfait, que la
quasi-projectivité de $X$ permet de représenter globalement par un complexe
borné $E^\bullet$ de Modules localement libres, en degrés $\geq 0$ ; le
théorème de comparaison (EGA~III, §~4.1) identifie $H^i(W_{/Y_W},
F_{/Y_W})$ à l'hypercohomologie de $\widehat{E}^\bullet$ ; et une récurrence
sur la longueur de $E^\bullet$, avec le lemme des cinq, fait passer
$(\mathrm{LG})_n$ des $E^p$ au complexe. Pour (ii), on représente
$U \mapsto f_* \underline{\mathrm{Hom}}(F, G \otimes f^* M)$ par un Module
cohérent (EGA~III, §~7.7), auquel on applique la bijectivité en degré $0$ et
l'injectivité en degré $1$ ; c'est pourquoi seule la platitude de $G$ compte.
Nous ne vérifions pas (iii).\footnote{Ces esquisses sont les nôtres. L'énoncé
(iii), et le Corollaire 2 qui en dépend, restent des énoncés du manuscrit.}

\emph{Corollaire 2} (page 42). Si $(X, Y)$ vérifie LG, le foncteur
$W \mapsto W \times_X X_{/Y}$, des schémas projectifs et plats sur $X$ dans
les schémas formels projectifs et plats sur $X_{/Y}$, est pleinement fidèle,
et de même en remplaçant $X$ par un voisinage ouvert de $Y$. Si $(X, Y)$
vérifie LG stricte, tout schéma formel projectif et plat sur $X_{/Y}$ est, sur
un voisinage ouvert $U$ de $Y$, le complété d'un schéma projectif et plat sur
$U$.

\emph{Corollaire 1} (page 42). Si $X$ est quasi-projectif et si $(X, Y)$
vérifie $(\mathrm{LG})_n$ (resp.\ LG, resp.\ LG stricte), alors pour tout
$W$ projectif et plat sur $X$, le couple $(W, f^{-1}(Y))$ vérifie
$(\mathrm{LG})_n$ (resp.\ LG, resp.\ LG stricte).\footnote{Le Corollaire 1 est
écrit après le Corollaire 2 sur la page ; « ($X$ quasi-projectif) » est
encadré. Pour $(\mathrm{LG})_n$, c'est (i) appliqué aux Modules localement
libres sur $W$ ; pour LG, on passe par $(\mathrm{LG})_0$ et la conclusion
pratique, $W$ étant quasi-projectif ; pour LG stricte, il faut (iii), et un
Module dont le complété le long de $f^{-1}(Y)$ est localement libre l'est au
voisinage de $f^{-1}(Y)$.}

\emph{Remarques} (page 43). Des exemples montrent que les conditions de
platitude et de propreté sont essentielles. Il est possible en revanche que
l'hypothèse de quasi-projectivité sur $X$ et $W$ puisse être omise --- ainsi
que, en partie, dans le Corollaire 2 (ii), l'« existence d'un schéma
algébrique donnant un schéma formel ».

\pagerange{45}{51}
\subsection*{Une forme locale (pages 45 à 51)}

\emph{Théorème.} Soient $A$ un anneau local noethérien, quotient d'un anneau
régulier, $t \in \mathfrak{m}$, et supposons $A$ séparé et complet pour la
topologie $t$-adique. Posons $Y = V(t)$, $Y' = Y \cap X'$,
$\widehat{X} = X_{/Y}$, $\widehat{X}' = X'_{/Y'}$. Les parties fermées de $X'$
disjointes de $Y'$ sont les parties finies formées de points fermés de $X'$ ;
notons $\Phi$ leur ensemble, $\mathcal{C}(X')$ la catégorie des Modules
cohérents sur $X'$ et $\mathcal{C}_\Phi(X')$ la sous-catégorie épaisse de ceux
dont le support est dans $\Phi$. Le foncteur
\[
  F \longmapsto \widehat{F} : \mathcal{C}(X') \longrightarrow \mathcal{C}(\widehat{X}')
\]
est exact, de noyau $\mathcal{C}_\Phi(X')$ ; il définit un foncteur exact et
fidèle
\[
  \rho : \mathcal{C}(X') / \mathcal{C}_\Phi(X') \longrightarrow \mathcal{C}(\widehat{X}') .
\]
Si $F$ et $G$ sont cohérents sur $X'$ et si $G$ est de profondeur $\geq 2$ aux
points fermés de $X' \smallsetminus Y'$, alors
\[
  \operatorname{Hom}(F, G) \xrightarrow{\ \sim\ } \operatorname{Hom}(\widehat{F}, \widehat{G}) .
\]
En particulier $\rho$ est pleinement fidèle sur la sous-catégorie pleine des
objets représentés par de tels Modules.\footnote{La page dit « $\rho$ est
pleinement fidèle », suivi d'une barre oblique, et parle d'une équivalence de
« $[\mathcal{C}(X')/\mathcal{C}_\Phi(X')]'$ » --- le prime marquant sans doute
une sous-catégorie --- avec les Modules cohérents sur $\widehat{X}'$ « qui sont
restriction de Modules cohérents sur $\widehat{X}$ ». Sans restriction, $\rho$
n'est pas plein : pour $A = \Bbbk[[x, y, z]]$, $t = x$ et $F = G =
\mathcal{O}_C$, $C = V(y)$, les endomorphismes de $F$ dans le quotient forment
l'anneau de valuation discrète $\Bbbk[[x, z]]_{(x)}$, ceux de $\widehat{F}$ son
complété $\Bbbk((z))[[x]]$. La complétude $t$-adique est fournie par nous : la
deuxième ligne de la page est surchargée, et on n'y lit que « la top.\
$t$-adique ». Elle est nécessaire, puisque $\widehat{X}'$ ne change pas quand
on remplace $A$ par son séparé complété $t$-adique, et que
$\Gamma(X', \mathcal{O})$ change.}

Par $H = \underline{\mathrm{Hom}}(F, G)$, qui est de profondeur
$\geq \min(2, \operatorname{prof} G)$ en chaque point, on se ramène à
$\Gamma(X', H) \simeq \Gamma(\widehat{X}', \widehat{H})$ pour $H$ de profondeur
$\geq 2$ aux points fermés de $X' \smallsetminus Y'$. La page 47 l'aborde par
un dévissage « facile » et « purement local » : une suite exacte
\[
  0 \to P \to H \to \bar{H} \to Q \to 0,
\]
où $P$ et $Q$ ont leur support dans $Y'$, de dimensions $\leq 0$ et
$\leq 1$, et $\bar{H}$ a la profondeur voulue ; puis, $K$ étant le noyau de la
multiplication par une puissance $t^n$ et $0 \to K \to H \to G_1 \to 0$ avec
$G_1$ sans $t$-torsion, la comparaison
\[
\begin{array}{ccccccc}
  0 \to \Gamma K & \longrightarrow & \Gamma H & \longrightarrow & \Gamma G_1 & \longrightarrow & H^1 K \\
  \downarrow \wr & & \downarrow & & \downarrow & & \downarrow \wr \\
  0 \to \Gamma \widehat{K} & \longrightarrow & \Gamma \widehat{H} & \longrightarrow & \Gamma \widehat{G}_1 & \longrightarrow & H^1 \widehat{K}
\end{array}
\]
où les flèches extrêmes sont des isomorphismes, $K$ étant tué par $t^n$. Le
cas $t$-régulier reste à faire : la page marque d'un « ? » la flèche de
$\Gamma G_1$, et s'arrête là.\footnote{La lettre lue $G'$ (ici $G_1$) pourrait
être $C'$ ; « $t^n$ » est incertain ; les bornes $\leq 0$, $\leq 1$ et la
barre de $\bar{F}'$ sont lues d'après le tracé, et la borne de profondeur de
$\bar{F}'$ est illisible. Le « ? » est de lui.}

Les pages 49 et 51 traitent l'autre moitié --- l'image essentielle. Avec un «
cor.~3 au th.\ de comparaison affine » et des conditions a), b), c') qui ne
sont pas dans le dossier --- c') disant que le Module formel $\mathcal{F}$ sur
$\widehat{X}'$ s'étend en un Module cohérent sur $\widehat{X}$ ---, on sait
que $\mathcal{F}$ provient d'un Module cohérent $F$ sans $t$-torsion, et l'on
veut un $F$ de profondeur $\geq 2$ aux points d'un ensemble $T_1$, qui est
sans doute celui des points fermés de $X' \smallsetminus Y'$. L'argument :
soit $U$ l'ouvert de $X'$ où cette profondeur est atteinte ; c'est un
voisinage de $Y'$, et ses points complémentaires sont fermés dans $X'$, donc
en nombre fini. Un point associé $x_1$ de $F$ dans $U$ dont l'adhérence
contiendrait un point $x$ de $X' \smallsetminus U$ serait d'adhérence de
dimension $1$, et $F$ serait de profondeur $\leq 1$ aux points fermés de cette
adhérence dans $U$, ce qui est absurde. Le critère de finitude des images
directes montre alors que, pour $i\colon U \to X'$ l'inclusion,
$F_1 = i_*(F|_U)$ est cohérent ; il est de profondeur $\geq 2$ aux points de
$X' \smallsetminus U$ par construction, et aux points fermés de $U$ par
définition de $U$ : « on a gagné ».\footnote{Le milieu de la page 49 est
encadré et annulé d'un long zigzag ; il définissait l'ensemble des points où
$\operatorname{prof} F_{x} \leq 1$. Une note marginale, « Je suppose $X$
localisé autour de $T$, et $X - T = X'$ », suggère un cadre plus général que
celui de la page 45, avec $T$ fermé. Nous suivons le texte non biffé.}

\section*{III. Morphismes projectifs : profondeur et dualité}

\pagerange{55}{58}
\subsection*{Les questions, et la dualité (pages 55 à 58)}

Soient $X = \mathbf{P}^r_{\Bbbk}$, $S = \Bbbk[T_0, \ldots, T_r]$, $Y$ une
section hyperplane, $U$ un voisinage ouvert de $Y$, $g\colon U \to X$
l'inclusion, $Z = X \smallsetminus U$, et $F$ un Module cohérent sur $X$. Le
fermé $Z$ est contenu dans l'espace affine $X \smallsetminus Y$ et fermé dans
$\mathbf{P}^r$ : c'est un ensemble \emph{fini} de points
fermés.\footnote{La page ne le dit pas ; c'est ce qui permet d'écrire, comme
elle le fait, la cohomologie à supports dans $Z$ comme un produit sur les
points $z \in Z$, et de parler de « $z$ isolé ».} Les questions :

a) À quelle condition $H^i(U, F(-m)) = 0$ pour $m$ grand ?

a') À quelle condition $\bigoplus_{m \geq 0} H^i(U, F(m))$ est-il de type fini
sur $S$ ?

b) À quelle condition $H^i(U, F(m)) = 0$ pour $m$ grand ($i > 0$ donné) ?

c) À quelle condition $H^i(U, F) \to H^i(\widehat{X}, \widehat{F})$ est-il
bijectif, injectif ?

(1°) On part de la suite exacte de cohomologie locale
\[
  \begin{array}{l}
  H^i_Z(X, F(-m)) \xrightarrow{\alpha_i} H^i(X, F(-m)) \to H^i(U, F(-m)) \\
  \qquad \to H^{i+1}_Z(X, F(-m)) \xrightarrow{\alpha_{i+1}} H^{i+1}(X, F(-m))
  \end{array}
\]
et l'on pose $E^j = \underline{\mathrm{Ext}}^j_{\mathcal{O}_X}(F, \omega_X)$.
Pour $m$ grand, la dualité de Serre identifie le dual de $H^i(X, F(-m))$ à
$\Gamma(X, E^{r-i}(m))$, et la dualité locale identifie le dual (de Matlis) de
$H^i_Z(X, F(-m)) = \bigoplus_{z \in Z} H^i_z(F_z)$ à
$\prod_{z \in Z} \widehat{E^{r-i}_z}$ ; l'application duale de $\alpha_i$ est
la restriction
\[
  \Gamma(X, E^{r-i}(m)) \longrightarrow \prod_{z \in Z} \widehat{E^{r-i}_z} .
\]
Comme la source est de dimension finie et que $\widehat{E^{r-i}_z}$ ne l'est
que si $E^{r-i}_z$ est de longueur finie, on obtient, pour $i$ donné :
\[
\begin{array}{l}
  \alpha_i \text{ injectif pour } m \text{ grand} \\
  \Updownarrow \\
  \text{tout } z \in Z \text{ est isolé dans } \operatorname{Supp} E^{r-i} \cup \lbrace z \rbrace \\
  \Updownarrow \\
  H^{i-1}(X, F(-m)) \to H^{i-1}(U, F(-m)) \text{ surjectif pour } m \text{ grand.}
\end{array}
\]
\footnote{« injectif » est écrit sur un premier mot, peut-être « surjectif » ;
la lecture « injectif » est celle que la dualité confirme. Les flèches duales
sont écrites sous $\alpha_i$ et $\alpha_{i+1}$. On utilise que $F(-m)$ et $F$
ont même cohomologie locale en un point, $\mathcal{O}(-m)$ y étant trivial.}

En prenant ces conditions pour tous les $i \leq n$, la dualité locale en un
point fermé $x$ ($E^{r-i}_x = 0$ si et seulement si $H^i_x(F_x) = 0$) les
traduit en profondeur :
\[
(*) \quad
\begin{array}{l}
  \alpha_i \text{ injectif pour } m \text{ grand},\ i \leq n \\
  \Updownarrow \\
  \exists\, V \text{ voisinage ouvert de } Z \text{ tel que } \operatorname{prof} F_x > n \text{ en tout point fermé } x \text{ de } U \cap V \\
  \Updownarrow \\
  H^i(X, F(-m)) \to H^i(U, F(-m)) \text{ surjectif pour } m \text{ grand},\ i < n .
\end{array}
\]
\footnote{La borne de profondeur n'est pas écrite sur la page ; nous la
fournissons. Elle entraîne, aux points non fermés $x$ de $U \cap V$,
$\operatorname{prof} F_x > n - \dim \overline{\lbrace x \rbrace}$, qui est la
condition $E^{r-i}_x = 0$ pour $i \leq n$.}

À la page 57, la même analyse pour la bijectivité. L'application duale est
bijective pour $m$ grand si et seulement si $E^{r-i}$ est à support dans
$Z$ : la finitude force les $E^{r-i}_z$ à être de longueur finie, puis
l'injectivité force la partie de $E^{r-i}$ à support hors de $Z$ à être nulle.
D'où, pour $i$ donné, puis pour tous les $i \leq n$ :
\[
\begin{array}{l}
  \alpha_i \text{ bijectif pour } m \text{ grand} \iff \operatorname{Supp} E^{r-i} \subset Z \\
  \qquad \iff H^i_x(F_x) = 0 \text{ pour tout } x \text{ fermé dans } U ; \\[1ex]
  \alpha_i \text{ bijectif pour } m \text{ grand},\ i \leq n
  \iff \operatorname{prof} F_x > n \text{ pour tout } x \text{ fermé dans } U \\
  \qquad \iff H^i(U, F(-m)) = 0 \text{ pour } i < n, \text{ et} \\
  \qquad\qquad H^n(U, F(-m)) \to H^{n+1}_Z(X, F(-m)) \text{ injectif}, \text{ pour } m \text{ grand.}
\end{array}
\]
\footnote{La dernière équivalence se lit sur la suite exacte : $\alpha_{i+1}$
injectif et $\alpha_i$ surjectif disent que $H^i(X) \to H^i(U)$ est à la fois
surjectif et nul. La page écrit à cet endroit une flèche
« $H^i(U, F(-m)) \Leftarrow H^i(X, F(-m))$ pour $m$ grand, $i < n$ », surchargée,
après avoir biffé « injectif pour $i \leq n$ » ; nous écrivons ce que donne la
suite exacte. Pour le passage d'un $x$ non fermé à un point fermé, voir la
note précédente.}

La condition qui répond à la question a) est la suivante, encerclée
$\mathrm{B}_n$ :
\[
(\mathrm{B}_n) \quad
\begin{array}{l}
  \alpha_i \text{ bijectif pour } m \text{ grand},\ i \leq n,
  \text{ et } \alpha_{n+1} \text{ injectif pour } m \text{ grand} \\
  \Updownarrow \\
  \operatorname{Supp} E^{r-i} \subset Z \text{ pour } i \leq n, \text{ et tout } z \in Z \text{ est isolé dans } \operatorname{Supp} E^{r-n-1} \cup \lbrace z \rbrace \\
  \Updownarrow \\
  \operatorname{prof} F_x > n \text{ pour } x \text{ fermé dans } U, \text{ et } \operatorname{prof} F_x > n+1 \text{ pour } x \text{ fermé dans } U \cap V, \\
  \quad V \text{ voisinage ouvert convenable de } Z \\
  \Updownarrow \\
  H^i(U, F(-m)) = 0 \text{ pour } m \text{ grand},\ i \leq n \\
  \Downarrow \\
  \bigoplus_{m \geq 0} H^i(U, F(m)) \text{ de type fini sur } S \text{ pour } i \leq n .
\end{array}
\]
Elle entraîne en particulier que les $H^i(U, F)$, $i \leq n$, sont de
dimension finie sur $\Bbbk$.\footnote{Deux écarts avec la page. (1) Elle
n'écrit qu'une implication $\Downarrow$ entre la deuxième et la troisième
ligne ; la réciproque vaut, par le calcul de la note sur $(*)$. (2) Elle écrit
une équivalence à la dernière ligne ; seule l'implication vaut. Si $F$ est le
faisceau gratte-ciel $\Bbbk$ en un point fermé de $U$, $\bigoplus_m H^0(U,
F(m)) \simeq \Bbbk[T]$ est de type fini sur $S$ mais $H^0(U, F(-m)) \neq 0$
pour tout $m$. Ce que la dernière ligne mesure est la condition
$\mathrm{A}^f_n$ ci-dessous, plus faible. La fin de la deuxième ligne est
illisible ; nous la complétons d'après (*). La conséquence de finitude est
écrite en marge, avec « cf 1° ».}

\emph{Le théorème de comparaison qui en découle.} Si $\mathrm{B}_n$ est
satisfaite et si l'équation $t$ de $Y$ est $F$-régulière, alors
\[
  H^i(U, F) \longrightarrow H^i(\widehat{X}, \widehat{F}) \simeq \varprojlim_m H^i(X_m, F_m)
\]
est bijectif pour $i < n$ et injectif pour $i = n$, $X_m$ désignant le
sous-schéma défini par $t^{m}$. On le voit sur les suites
$H^i(U, F(-m)) \to H^i(U, F) \to H^i(X_m, F_m) \to H^{i+1}(U, F(-m))$ déduites
de $0 \to F(-m) \xrightarrow{t^m} F \to F_m \to 0$.\footnote{La page écrit «
sous les conditions de $(*)$ » ; ce sont les annulations de $\mathrm{B}_n$, que
la page vient d'établir (« Cela résout la question (a) »), qui servent, et les
conditions $(*)$ ne les donnent pas. L'isomorphisme
$H^i(\widehat{X}, \widehat{F}) \simeq \varprojlim H^i(X_m, F_m)$ vaut ici
pour tout $i$, les $H^i(X_m, F_m)$ étant de dimension finie.}

Ces conditions suffisent pour c), mais ne sont pas nécessaires : pour $n = 1$,
il suffit que $\operatorname{prof} F_x \geq 2$ aux points fermés de $U$ pour
que $H^0(U, F) \to H^0(\widehat{X}, \widehat{F})$ soit
bijectif.\footnote{L'énoncé est de la page 58. On peut le voir ainsi :
$g_*(F|_U)$ est alors cohérent, de profondeur $\geq 2$ en tout point fermé de
$X \smallsetminus Y$, et le théorème de comparaison de SGA~2, XII, pour $X$
projectif s'applique ; la justification est la nôtre.}

\emph{Deux brouillons} (pages 56 et 60). Au dos de deux exemplaires d'une
même page dactylographiée, qui s'achève sur l'annonce d'un « Corollaire 3.2 »,
Grothendieck calcule à l'encre puis au crayon : des sommes
$\coprod_{p, q \geq 0} F_0(q - p)$ et
$\sum_{p \geq 0} R^i f_{0*} F_0(q - p)$ pour $q - p \geq n$, un module gradué
$\coprod_q S_q$, un double indice $S_{pq}$ nul pour $p > 0$ et égal à $S_q$
pour $p = 0$, et des quadrants hachurés en $(p, q)$ légendés
« $\mathcal{F}(n)$ » et « filtre ». On y reconnaît le calcul d'une
filtration indexée par deux degrés, sans pouvoir le rattacher sûrement à une
question du dossier.\footnote{Le texte dactylographié porte sur l'existence
d'un Module cohérent de complété formel donné, sous des hypothèses de
régularité, de platitude et de profondeur $\geq 2$ ; on n'en dit pas plus que
la transcription. L'indice $pq$ de $S$ est d'une lecture incertaine.}

\pagerange{58}{61}
\subsection*{Finitude, et torsions positives (pages 58 à 61)}

(2°) La suite exacte montre que $H^i(U, F) \to H^{i+1}_Z(X, F)$ est un
isomorphisme modulo des espaces vectoriels de dimension finie. D'où :
\[
(\mathrm{A}'_i) \quad
\begin{array}{l}
  H^i(U, F) \text{ de dimension finie sur } \Bbbk \\
  \Updownarrow \\
  H^{i+1}_z(F_z) \text{ de longueur finie pour tout } z \in Z \\
  \Updownarrow \\
  \text{tout } z \in Z \text{ est isolé dans } \operatorname{Supp} E^{r-i-1} \cup \lbrace z \rbrace \\
  \Updownarrow \\
  \alpha_{i+1} \text{ injectif pour } m \text{ grand} \\
  \Updownarrow \\
  H^i(X, F(-m)) \to H^i(U, F(-m)) \text{ surjectif pour } m \text{ grand.}
\end{array}
\]

(3°) Pour $i \geq 1$, $R^i g_*(F|_U) = \underline{H}^{i+1}_Z(F)$, et
$g_*(F|_U)$ est cohérent si et seulement si $\underline{H}^1_Z(F)$ l'est ; un
faisceau à support dans l'ensemble fini $Z$ est cohérent si et seulement si
ses fibres sont de longueur finie. Pour $m$ grand, $H^i(U, F(m))$ s'identifie à
$H^{i+1}_Z(X, F)$ ($i \geq 1$), et la multiplication par l'équation de $Y$,
inversible au voisinage de $Z$, y est bijective. On en tire :
\[
(\mathrm{A}^f_n) \quad
\begin{array}{l}
  H^i(U, F) \text{ de dimension finie pour } i \leq n \\
  \Updownarrow \\
  \operatorname{prof} F_x > n+1 \text{ aux points fermés } x \text{ de } U \text{ voisins de } Z \\
  \Updownarrow \\
  \text{tout } z \in Z \text{ est isolé dans } \operatorname{Supp} E^{r-i-1} \cup \lbrace z \rbrace,\ i \leq n \\
  \Updownarrow \\
  \alpha_i \text{ injectif pour } m \text{ grand},\ i \leq n+1 \\
  \Updownarrow \\
  H^i(X, F(-m)) \to H^i(U, F(-m)) \text{ surjectif pour } m \text{ grand},\ i \leq n \\
  \Updownarrow \\
  R^i g_*(F|_U) \text{ cohérent pour } i \leq n
  \iff H^i_z(F_z) \text{ de longueur finie pour } z \in Z,\ i \leq n+1 \\
  \Updownarrow \\
  \bigoplus_{m \geq 0} H^i(U, F(m)) \text{ de type fini sur } S \text{ pour } i \leq n .
\end{array}
\]
\footnote{Les labels $\mathrm{A}'_i$ et $\mathrm{A}^f_n$ sont encerclés en
marge ; l'exposant $f$ est incertain. Une double flèche de la page n'a qu'une
pointe. La deuxième ligne est, par le théorème de finitude de SGA~2, VIII, la
condition que les $H^j_z(F_z)$, $j \leq n+1$, soient de longueur finie. Le haut
de la page 59, annulé, portait $R^p f_* [R^q j_* (F(m))]$ : la suite
spectrale de Leray par laquelle on passe de $R^i g_*$ aux modules gradués.}

Les questions a) et b) portent ainsi sur des points différents : a) sur la
profondeur aux points fermés de $U$, b) --- les torsions positives --- sur les
points de $Z$. La page 61 le dit :
\[
(\mathrm{C}'_i) \quad
\begin{array}{l}
  H^i(U, F(m)) = 0 \text{ pour } m \text{ grand } (i > 0 \text{ donné}) \\
  \Updownarrow \\
  \underline{H}^{i+1}_Z(F) = 0, \text{ i.e. } H^{i+1}_z(F_z) = 0 \text{ pour } z \in Z \\
  \Updownarrow \\
  R^i g_*(F|_U) = 0 ;
\end{array}
\]
\[
(\mathrm{C}_n) \quad
\begin{array}{l}
  H^i(U, F(m)) = 0 \text{ pour } m \text{ grand},\ 0 < i \leq n \\
  \Updownarrow \quad (\text{si } F = g_*(F|_U)) \\
  \operatorname{prof}_z F > n+1 \text{ pour } z \in Z .
\end{array}
\]
« Les conditions $C$ sont les seules qui exigent quelque chose sur la
profondeur en les $z \in Z = X - U$\ldots\ et non seulement aux points de $U$
! »\footnote{La page écrit $\Downarrow$ puis une double flèche surchargée entre
les deux dernières lignes de $\mathrm{C}'_i$ ; l'équivalence vaut, puisque
$R^i g_*(F|_U) = \underline{H}^{i+1}_Z(F)$ pour $i \geq 1$. La lecture
« supposant que » devant $F = g_*(F|_U)$ est incertaine ; cette hypothèse
annule $H^0_z$ et $H^1_z$, et c'est elle qui rend la seconde équivalence
exacte. Le texte s'arrête sur un trait ondulé.}

\section*{IV. Le groupe de Picard des cônes}

\pagerange{64}{66}
\subsection*{Brouillons (pages 64 à 66)}

Les pages 64 et 65 sont deux feuilles de brouillon, faites de diagrammes, de
formules et de petites figures, et la page 66 un départ au propre qui
s'interrompt. Le sujet est le point 9 du plan : comparer le groupe de Picard
d'une variété projective à celui de son cône, ce que demande un théorème de
Lefschetz pour $\operatorname{Pic}$ dans le cas local --- le cône affine en
son sommet étant l'exemple type d'anneau local.

\emph{La situation de la page 66.} Soient $X$ un schéma projectif sur
$\Bbbk$ avec $\Gamma(X, \mathcal{O}_X) = \Bbbk$, $L = \mathcal{O}_X(1)$ ample,
$R = \bigoplus_{m \geq 0} \Gamma(X, L^{\otimes m})$, $C = \operatorname{Spec}
R$ le cône projetant affine, de sommet $a$, et $\widehat{C} =
\operatorname{Proj} R[z]$ le cône projetant projectif.\footnote{La page 64
écrit « $\widehat{C} = \operatorname{Spec} S[z]$ », le mot « Spec » d'une
lecture incertaine ; pour le cône projectif, c'est $\operatorname{Proj}$. Le
renvoi de la page 66 à « II 8.7.5 » et « II 8.8.6 » est, sans doute, à EGA~II,
§~8, sur les cônes projetants ; il n'a pas été vérifié. Prendre pour $R$
l'anneau des sections, plutôt qu'un anneau de coordonnées quelconque, est ce
qui donne $f_* \mathcal{O}_P = \mathcal{O}_{\widehat{C}}$.} Soit $P$ l'éclaté de
$\widehat{C}$ en son sommet : c'est le fibré en droites projectif
$\pi\colon P \to X$, muni de la section $i\colon X \to P$ (le diviseur
exceptionnel) et de la contraction $f\colon P \to \widehat{C}$ qui envoie
$i(X)$ sur $a$.\footnote{La page note $X'$ ce que nous notons $P$.}

\begin{tikzcd}[column sep=small, row sep=small, nodes={font=\scriptsize}]
X \arrow[r, "i"] \arrow[d] & P \arrow[r, "\pi"] \arrow[d, "f"] & X \\
a \arrow[r] & \widehat{C} &
\end{tikzcd}

Alors :
\begin{itemize}
\item $\operatorname{Pic}(P) \simeq \operatorname{Pic}(X) \times \mathbf{Z}$
(fibré projectif), et le noyau de
$i^*\colon \operatorname{Pic}(P) \to \operatorname{Pic}(X)$ est isomorphe à
$\mathbf{Z}$, engendré par $\mathcal{O}_P(1) = f^* \mathcal{O}_{\widehat{C}}(1)$ ;
\item $\operatorname{Pic}(\widehat{C}) \to \operatorname{Pic}(P)$ est injectif,
car $f_* \mathcal{O}_P = \mathcal{O}_{\widehat{C}}$ ;
\item d'où, comme il le note encadré (« je m'aperçois que »),
\[
  \operatorname{Pic}(\widehat{C}) \simeq \mathbf{Z} \xrightarrow{\ \sim\ }
  \operatorname{Ker}\bigl[\operatorname{Pic}(P) \xrightarrow{i^*} \operatorname{Pic}(X)\bigr] .
\]
\end{itemize}

Soit enfin $C_X = P \smallsetminus f^{-1}(\widehat{C} \smallsetminus C)$ la
partie de $P$ au-dessus de $C$ : c'est l'espace total d'un fibré en droites sur
$X$, l'éclaté de $C$ en son sommet. La page affirme, avec un point
d'interrogation, que $\pi^*\colon \operatorname{Pic}(X) \to
\operatorname{Pic}(C_X)$ est un isomorphisme, et demande en marge s'il faut
supposer $X$ localement factoriel. L'injectivité est démontrée : le composé
avec la restriction à la section nulle est l'identité. La démonstration de la
surjectivité commence --- « au-dessus d'un ouvert convenable de $X$, » --- et
s'arrête là, au milieu de la phrase : c'est la dernière ligne du
dossier.\footnote{Réponse à la question de la marge, qui est la nôtre : la
factorialité locale n'est pas nécessaire, mais une hypothèse l'est. Si $X$ est
réduit et seminormal --- normal par exemple ---, $\operatorname{Pic}$ est
invariant par la droite affine sur les anneaux locaux de $X$ (théorème de
Traverso et Swan, 1970-1980), et $\operatorname{Pic}(C_X) = \operatorname{Pic}(X)$.
Pour une cubique cuspidale, il est faux : $\operatorname{Pic}(A[T]) \neq
\operatorname{Pic}(A)$ pour $A = \Bbbk[t^2, t^3]$.}

Les brouillons des pages 64 et 65 entourent cette question. On y lit :
\begin{itemize}
\item une suite $\operatorname{Pic}(X) \to \operatorname{Pic}(\widehat{C}') \to
\operatorname{Pic}(C') \to \operatorname{Pic}(\widetilde{C'})$, annotée
« injectif », « bijectif dans le cas $X$ loc.\ factoriel », un noyau
« engendré par l'image de $\mathcal{O}_X(1)$ », et les quotients
$\operatorname{Pic}(X)/\mathbf{Z}\xi \subset
\operatorname{Pic}(\widehat{C}')/\mathbf{Z}\xi$ ;\footnote{Les objets
primés et tildés ne sont définis nulle part ; les étiquettes « injectif »,
« surjectif » sont mal rattachées à leurs flèches. Nous ne les identifions
pas.}
\item deux fois, encadrée, la question $\operatorname{Pic}(C) = 0$ ?? ---
qui a une réponse : oui si $R$ est seminormal (normal par exemple), car alors
$\operatorname{Pic}(R) = \operatorname{Pic}(R[T])$, et l'homomorphisme
$R \to R[T]$, $r \mapsto T^{\deg r} r$, relie l'identité ($T = 1$) à la
projection sur $R_0 = \Bbbk$ ($T = 0$) ; la réponse peut être non pour un $R$ qui
ne l'est pas, $\Bbbk[t^2, t^3]$ en étant l'exemple type ;\footnote{La réponse est la nôtre.}
\item la suite de cohomologie locale de $\mathcal{O}_X^*$ pour
$U = X \smallsetminus \lbrace a \rbrace$,
\[
  \begin{array}{l}
  0 \to H^0_a(\mathcal{O}^*) \to H^0(X, \mathcal{O}^*) \to H^0(U, \mathcal{O}^*)
  \to H^1_a(\mathcal{O}^*) \to H^1(X, \mathcal{O}^*) \\
  \qquad \to H^1(U, \mathcal{O}^*) \to H^2_a(\mathcal{O}^*) \to H^2(X, \mathcal{O}^*),
  \end{array}
\]
avec $H^1(X, \mathcal{O}^*) = 0$ pour $X$ local --- de sorte que
$\operatorname{Pic}(U)$ et le défaut de prolongement des unités se lisent dans
$H^1_a$ et $H^2_a$ : c'est la notion d'anneau \emph{parafactoriel}
(EGA~IV$_4$, §~21.13), pour lequel $\operatorname{Pic}(U) = 0$ et
$\Gamma(U, \mathcal{O}) = \Gamma(X, \mathcal{O})$ ;
\item la hiérarchie « loc.\ factoriel $\Rightarrow$ loc.\ quasi-factoriel
$\Rightarrow$ parafactoriel » et « factoriel $\Rightarrow$ $1$-factoriel
$\Rightarrow$ $0$-factoriel » ;\footnote{Les notions intermédiaires ne sont
pas définies sur ces pages. Un anneau local factoriel n'est parafactoriel
qu'en dimension $\geq 2$.}
\item des formules de genre pour une courbe à points doubles : si $C$ a $\nu$
composantes de genres $g_i$ et des points singuliers $s$ de contributions
$n_s$, $1 - g = \sum (1 - g_i) - \sum n_s$, soit
$g = 1 + \sum g_i + \sum n_s - \nu$ ; avec les deux cas $g = g' + \sum n_s$
(une composante) et $g = g' + g'' + \sum n_s - 1$ (deux), et l'exemple
$\Bbbk[x, y, z]/(xy)$.\footnote{Ces formules sont celles du genre
arithmétique, $n_s$ étant l'invariant $\delta$ du point $s$ (le nombre de
points doubles ordinaires). La formule « $1 - h' = 2 - h' - h''$ », dont le
premier terme est incertain, et les relations $g^m = u f^n$, $g = u f^k$ de la
page 64 ne sont pas rattachées.}
\end{itemize}

\end{document}
